% Содержимое подраздела 3.1 - Стандартизация оценки: бенчмарки и метрики


\section{Стандартизация оценки: бенчмарки и метрики}
\label{sec:standardization}

Одной из центральных проблем, тормозивших развитие области генеративных моделей, было отсутствие единого стандарта для их сравнения. Исследователи использовали различные наборы данных, методы их предобработки и метрики качества, что делало невозможным объективное сравнение эффективности новых архитектур по принципу <<яблоки с яблоками>>. Это приводило к несопоставимым, зачастую невоспроизводимым результатам и затрудняло оценку реального прогресса \cite{polykovskiy2020molecular}.

Ответом на этот вызов стала разработка комплексной платформы с открытым исходным кодом MOSES (Molecular Sets) \cite{polykovskiy2020molecular}. Данная платформа не просто предложила новый метод, а создала инфраструктуру для стандартизированного обучения и оценки, состоящую из трех ключевых компонентов:
\begin{enumerate}
    \item \textbf{Стандартизированный набор данных.} Авторы MOSES тщательно отфильтровали коллекцию ZINC Clean Leads, создав датасет с четко определенным разделением на обучающую, тестовую и, что особенно важно, \textit{scaffold test} выборки. Последняя содержит молекулы с структурными каркасами, отсутствующими в обучающей выборке, что позволяет оценить способность модели к генерализации и созданию принципиально новых химических структур.
    \item \textbf{Реализация базовых моделей.} Платформа включает эталонные реализации нескольких ключевых архитектур (CharRNN, VAE, AAE, JTN-VAE), которые служат <<точкой отсчета>> (reference point) для всех будущих моделей.
    \item \textbf{Комплексный набор метрик.} MOSES формализовал и популяризовал набор метрик для всесторонней оценки генеративных моделей, который впоследствии стал стандартом де-факто в большинстве последующих работ \cite{mouchlis2021advances, kim2021comprehensive}.
\end{enumerate}

Ключевые метрики, предложенные в MOSES, можно разделить на несколько категорий. Базовые метрики качества генерации включают:
\begin{itemize}
    \item \textbf{Validity (Валидность):} Процент сгенерированных строк (например, SMILES), которые являются синтаксически и семантически корректными и могут быть преобразованы в валидную химическую структуру.
    \item \textbf{Uniqueness (Уникальность):} Процент уникальных молекул среди всех валидных сгенерированных. Низкое значение этого показателя может свидетельствовать о коллапсе мод (mode collapse), когда модель генерирует лишь небольшое подмножество известных ей структур.
    \item \textbf{Novelty (Новизна):} Процент уникальных валидных молекул, отсутствующих в обучающем наборе данных. Эта метрика напрямую измеряет способность модели создавать новые, ранее не виденные соединения.
    \item \textbf{Internal Diversity (IntDiv, Внутреннее разнообразие):} Среднее попарное несходство (обычно 1 - Tanimoto similarity) молекул внутри сгенерированного набора. Высокое значение указывает на то, что модель генерирует химически разнообразные структуры.
\end{itemize}

Однако наиболее важным вкладом MOSES стала популяризация метрик, оценивающих сходство распределений. Среди них центральное место занимает \textbf{Fréchet ChemNet Distance (FCD) \cite{polykovskiy2020molecular}. FCD измеряет расстояние Фреше между распределениями активаций предпоследнего слоя нейросети ChemNet (предобученной на предсказание биологической активности) для сгенерированного и референсного наборов молекул. Низкое значение FCD указывает на то, что сгенерированные молекулы не только структурно, но и с точки зрения потенциальных биологических свойств похожи на реальные данные. Эта метрика стала мощным интегральным показателем, позволяющим оценить, насколько хорошо генеративная модель уловила сложное многомерное распределение химически осмысленных структур.
